\documentclass[11pt]{article}

% ====================================================
% PACKAGES
% ====================================================
\usepackage[margin=1in, headheight=13.6pt]{geometry}
\usepackage[T1]{fontenc}
\usepackage{lmodern}
\usepackage{microtype}
\usepackage{setspace}
\usepackage{tcolorbox}
\usepackage{amsmath, amssymb, mathtools}
\usepackage{siunitx}
\usepackage{bm}

\usepackage{graphicx}
\usepackage{float}
\usepackage{subcaption}

\usepackage{enumitem}
\usepackage{booktabs}
\usepackage{longtable}
\usepackage{array}
\usepackage{multirow}

\usepackage{xcolor}
\usepackage{hyperref}

\usepackage{fancyhdr}
\usepackage{lastpage}

\usepackage{indentfirst}
\usepackage[justification=centering]{caption}
\usepackage{pdfpages}
\usepackage{listings}

\lstdefinestyle{pystyle}{
    language=Python,
    basicstyle=\ttfamily\small,
    keywordstyle=\color{blue!60!black}\bfseries,
    commentstyle=\color{green!40!black}\itshape,
    stringstyle=\color{orange!70!black},
    numbers=left,
    numberstyle=\tiny\color{gray},
    numbersep=6pt,
    breaklines=true,
    breakatwhitespace=false,
    showstringspaces=false,
    frame=single,
    rulecolor=\color{gray!50},
    tabsize=4
}
\lstset{style=pystyle}

\pagestyle{fancy}
\fancyhf{}
\lhead{\Assignment\ |\ \course}
\rhead{\Name\ |\ \today}
\cfoot{Page \thepage\ of~\pageref{LastPage}}

\hypersetup{
    colorlinks,
    linkcolor={red!50!black},
    citecolor={blue!50!black},
    urlcolor={blue!80!black}
}

\newcommand{\Name}{Arael Anaya}
\newcommand{\Assignment}{Homework 6}
\newcommand{\course}{Robot Mechanics}

% Boxed final-answer environment
\newtcolorbox{result}{
    colback=blue!3!white,
    colframe=blue!40!black,
    boxrule=0.6pt,
    arc=2pt,
    left=6pt, right=6pt, top=4pt, bottom=4pt
}

% Boxed problem-statement environment
\newtcolorbox{statement}{
    colback=gray!5!white,
    colframe=gray!50!black,
    boxrule=0.4pt,
    arc=2pt,
    left=6pt, right=6pt, top=4pt, bottom=4pt
}

% Visible placeholder for work not yet added
\newcommand{\workplaceholder}{%
    \par\vspace{4pt}\noindent\textcolor{gray!60}{\itshape [Work not yet added: paste solution here]}\par\vspace{4pt}%
}

% Boxed correction-note environment (flags a fix made to the original hand calc)
\newtcolorbox{fixnote}{
    colback=orange!5!white,
    colframe=orange!55!black,
    boxrule=0.5pt,
    arc=2pt,
    left=6pt, right=6pt, top=4pt, bottom=4pt,
    fontupper=\small
}

% Notation macros (matching course conventions)
\newcommand{\T}{\bm{T}}
\newcommand{\R}{\bm{R}}
\newcommand{\dvec}{\bm{d}}
\newcommand{\framen}[1]{\{#1\}}
\newcommand{\khat}{\widehat{\bm{k}}}
\newcommand{\wvec}{\bm{\omega}}
\newcommand{\xivec}{\bm{\xi}}
\newcommand{\vvec}{\bm{v}}
\newcommand{\Omvec}{\bm{\Omega}}
\newcommand{\Rotx}{\operatorname{Rot}_x}
\newcommand{\Roty}{\operatorname{Rot}_y}
\newcommand{\Rotz}{\operatorname{Rot}_z}
\newcommand{\atantwo}{\operatorname{atan2}}
\newcommand{\skewx}[1]{\left[#1\right]_{\times}}

\begin{document}



% ====================================================
% PROBLEM 1
% ====================================================
\section*{Problem 1}

\begin{statement}
A manipulator with one prismatic and five rotational joints has link lengths $a_1 = 1$ and $a_2 = 0.5$ and the DH table below ($\star$ marks a joint variable).
\begin{center}
\begin{tabular}{c|cccc}
Link & $a$ & $d$ & $\alpha$ & $\theta$ \\ \hline
1 & $a_1$ & 0 & 0 & $\star$ \\
2 & $a_2$ & 0 & 0 & $\star$ \\
3 & 0 & $\star$ & 0 & 0 \\
4 & 0 & 0 & $\pi/2$ & $\star$ \\
5 & 0 & 0 & $-\pi/2$ & $\star$ \\
6 & 0 & 0 & 0 & $\star$
\end{tabular}
\end{center}
The initial pose at $t=0$\,s and final pose at $t=2$\,s are
\[
{}^0\T_6(0)=\begin{bmatrix}1&0&0&1.00\\0&0&-1&0.25\\0&1&0&0.25\\0&0&0&1\end{bmatrix},\qquad
{}^0\T_6(2)=\begin{bmatrix}0&-1&0&0.50\\1&0&0&-0.50\\0&0&1&0\\0&0&0&1\end{bmatrix}.
\]
\begin{enumerate}[label=(\alph*), nosep]
\item Write a MATLAB function to generate the coefficients for cubic interpolation between two points given zero initial and final velocity.
\item Write a MATLAB function to interpolate between two vectors, given cubic interpolation coefficients, an initial time, and a list of times for evaluation.
\item For each parameterization below, generate a figure of the path in joint variable space and a figure of the Cartesian path, and play the motion as a movie using \texttt{drawArm.m}. Discuss the differences and similarities.
\begin{enumerate}[label=\roman*., nosep]
\item Parameter space $(\theta, d)$ from initial joint values to final.
\item Cartesian/angle-axis space.
\item Cartesian/quaternion space (normalize the interpolated quaternion at each sub point).
\item Twist space.
\end{enumerate}
\end{enumerate}
\end{statement}

% ----------------------------------------------------
\subsection*{Parts (a) and (b)}

\texttt{cubicCoeffs} returns the rest-to-rest cubic $q(t)=c_0+c_1\tau+c_2\tau^2+c_3\tau^3$, $\tau=t-t_0$, with $c_0=q_0$, $c_1=0$, $c_2=3(q_f-q_0)/T^2$, $c_3=-2(q_f-q_0)/T^3$, one row per element of $q$. \texttt{cubicInterp} evaluates these at a list of times, returning an $n\times m$ matrix. Code is in Appendix~A.

\subsection*{Part (c)}

The inverse kinematics of the supplied closed-form solution is implemented in \texttt{ik\_hw6}. The final pose has $r_{13}=r_{23}=0$ and $r_{33}=1$, so $\theta_5=0$ and the pose is a wrist singularity. In that case \texttt{ik\_hw6} sets $\theta_6=0$ and solves $\theta_4 = \operatorname{atan2}(r_{21}, r_{11}) - \theta_1 - \theta_2$. Each trajectory is interpolated in its own space with a rest-to-rest cubic over $0\le t\le2$\,s (50 samples), converted back to a pose at every sample, and passed through the inverse kinematics. End-effector positions for the Cartesian figures come from forward kinematics (\texttt{armPositions}). All four trajectories start and end at the specified poses to machine precision.

\newcommand{\trajfigs}[2]{%
    \begin{figure}[H]
        \centering
        \begin{subfigure}[t]{0.49\textwidth}
            \centering
            \includegraphics[width=0.85\linewidth]{Figures/P1_joint_#1.png}
            \caption{Joint variables vs.\ time}
        \end{subfigure}\hfill
        \begin{subfigure}[t]{0.49\textwidth}
            \centering
            \includegraphics[width=0.85\linewidth]{Figures/P1_cartesian_#1.png}
            \caption{Cartesian end-effector path}
        \end{subfigure}
        \caption{#2}
    \end{figure}%
}

\subsubsection*{i. Parameter space}
\trajfigs{jointspace}{Joint-space interpolation.}

\subsubsection*{ii. Cartesian/angle-axis space}
\trajfigs{angleaxis}{Cartesian/angle-axis interpolation.}

\subsubsection*{iii. Cartesian/quaternion space}
\trajfigs{quaternion}{Cartesian/quaternion interpolation.}

\subsubsection*{iv. Twist space}
\trajfigs{twist}{Twist-space interpolation.}

\subsubsection*{Discussion}

\begin{result}
Analyzing the four different interpolation methods, I found it interesting to note the difference in curvature over time between them. Angle-axis and quaternion both created a straight line from the start to the end pose. I believe this is because position is interpolated linearly in Cartesian space, independent of the rotation. We see that joint space and twist both create a much smoother, curved path that gradually gets to the end. However, only the joint-space method also has smooth joint curves, as is clear from the joint variables vs.\ time graphs I created. In parameter space, the joints have much smoother curves, implying a much smoother velocity and acceleration profile. The angle-axis, quaternion, and twist methods all show an abrupt jump in $\theta_4$ and $\theta_6$ in the final sample, because the final pose is a wrist singularity ($\theta_5 = 0$).
\end{result}

\bigskip
% ====================================================
% APPENDIX A: CODE
% ====================================================
\newpage
\section*{Appendix A: Code}

% Listings are skipped until the code file exists, so a missing file cannot break the build.
\newcommand{\codelisting}[2]{%
    \IfFileExists{../Code/#1}{\lstinputlisting[language=#2, caption={\detokenize{#1}}]{"../Code/#1"}}%
    {\noindent\textcolor{gray!60}{\itshape [#1 not yet written]}}%
}
\newcommand{\funclisting}[2]{%
    \IfFileExists{../../../Functions/#1}{\lstinputlisting[language=#2, caption={#1}]{"../../../Functions/#1"}}%
    {\noindent\textcolor{gray!60}{\itshape [#1 not yet written]}}%
}

\subsection*{cubicCoeffs.m (Part a)}\label{app:code:cubicCoeffs}
\codelisting{cubicCoeffs.m}{Matlab}

\subsection*{cubicInterp.m (Part b)}\label{app:code:cubicInterp}
\codelisting{cubicInterp.m}{Matlab}

\subsection*{ik\_hw6.m}\label{app:code:ik}
\codelisting{ik_hw6.m}{Matlab}

\subsection*{armPositions.m}\label{app:code:armPositions}
\codelisting{armPositions.m}{Matlab}

\subsection*{plotTrajectory.m}\label{app:code:plotTrajectory}
\codelisting{plotTrajectory.m}{Matlab}

\subsection*{Problem1.m (Part c)}\label{app:code:problem1}
\codelisting{Problem1.m}{Matlab}

% ====================================================
% APPENDIX B: HAND CALCULATIONS
% ====================================================
\newpage
% \section*{Appendix B: Hand Calculations}\label{app:hand}

% \IfFileExists{../Homework_6_hand_calcs.pdf}{%
%     \includepdf[pages=-, pagecommand={\thispagestyle{plain}}, scale=0.9]{../Homework_6_hand_calcs.pdf}%
% }{\noindent\textcolor{gray!60}{\itshape [Homework\_6\_hand\_calcs.pdf not yet added]}}

\end{document}
